\chapter{Drive Mechanisms}

Drive mechanisms are an important part of any automatic mechanical system as they are responsible for the transformation of motor action into linear or rotary movement. In gantry systems, drive mechanisms determine not only how motion is provided, but also influence the accuracy of positioning, load applied to the structure, efficiency of the system, as well as the reliability of its operation. Consequently, choosing an appropriate drive involves numerous factors, which include the carrying capacity, accuracy, speed, complexity of design, installation, and maintenance.

The choice of a drive mechanism is especially crucial when designing industrial material handling systems such as automated packaging or palletizing machines as they require efficient transmission of movement for relatively long distances and precise positioning. Moreover, the actuator system has to withstand constantly changing dynamic loads that arise due to acceleration, deceleration, and transfer of cargo. Hence, choosing a proper drive directly affects the operation and lifespan of the machinery.

This chapter covers the drive systems that have been studied in relation to the gantry system, especially the belt pulley and rack \& pinion drive system for horizontal and vertical movements.

\subsubsection{X-Axis Actuation Design}
For actuation of the X-axis, a belt and pulley mechanism system was selected over Rack and Pinion. The details of a belt and pulley mechanism are given below:

Belt and pulley systems are among the most commonly used transmission mechanisms in industrial machinery. These kind of  systems transmit mechanical power between rotating shafts using a flexible belt that runs over pulleys mounted on the shafts. The basic working principle involves transferring motion from the driving pulley, which is connected to a motor, to the driven pulley through the tensile forces developed in the belt. Due to their simplicity and versatility, belt drives are commonly used in machines requiring smooth motion transmission over moderate to long distances.  

The advantages of using the belt and pulley system are:
\begin{itemize}
    \item They are ideal for motion systems that require extended travel distances along an axis with precision
    \item The driving motor can remain stationary while the belt transmits motion along the axis. This reduces the mass of moving components and increase accelaration capabilty.
\item Belts can be installed and replacd relatively easily
\item Belts are also generally less expensive in terms of component cost and manufacturing complexity.

The system consist of timing belts routed over the  pulleys across the axis and the pulleys have long shaft connected to the one of the pulleys at one end for better motion. Also, a dual motor configuration was considered for the X-axis, allowing symmetric actuation across both sides of the structure.
\end{itemize}

\subsubsection{Z-Axis Lifting Mechanism}
The Z-Axis provides vertical lifting capability of the for the end-effector and the payload.

A rack and pinion mechanism is selected for the vertical actuation due to it's robustness. Rack-and-pinion systems also better resistance to slippage.
when lifting loads vertically.

During the design process, both helical and spur gear configurations were evaluated. However, spur gear were ultimately selected because they provide more reliable engagement and alignment for the vertical drive systems due to their straight teeths parallel to axis.

The motor is mounted with the motor housing, and couplers are used to connect the motor shaft, and the couplers are used to connect the motor shaft to the pinion gear which is driving the rack.

Additionally, to improve the stability and prevent lateral movements during lifting, linear guide channels were incorporated along the z-axis. Custom mounts were designed and fabricated using 3d printing to securely attached the linear guides and maintain allignment during operations.
